\documentclass{article}
\usepackage[utf8]{inputenc}
\usepackage[cm]{fullpage}
%\usepackage[spanish,mexico]{babel}

\usepackage{siunitx}
\sisetup{
	output-complex-root = \ensuremath{\mathrm{j}},
	complex-root-position = before-number
}

\usepackage[inline]{enumitem}

\usepackage{graphicx}
%\usepackage{wrapfig}
\usepackage{hyperref}
\usepackage{datatool}



\usepackage{tabularx}
\usepackage[table]{xcolor}
\usepackage{multirow}
\usepackage{hhline}


\definecolor{Gray}{gray}{0.85}
\definecolor{LightCyan}{rgb}{0.88,1,1}

\newcommand{\Subject}{Semiconductor physics}
\newcommand{\Group}{4A}
\newcommand{\Carrera}{Electrical engineering}
\newcommand{\ExamType}{Exam 2}
\newcommand{\Date}{15/05/2017}
\newcommand{\MaximumMarks}{10}
\newcommand{\PName}{Dr. Suresh Kumar Gadi}


\makeatletter
% Blank/missing fields commands
% \skipblank adds \\ to filled field; * version adds \space instead of newline
\newcommand\skipblank{\@ifstar\@spskip\@nlskip}
\newcommand\@nlskip[1]{\ifthenelse{\DTLiseq{#1}{}}{\relax}{#1\\}}
\newcommand\@spskip[1]{\ifthenelse{\DTLiseq{#1}{}}{\relax}{#1\space}}
% \checkblank replaces blank fields with ***
\newcommand\checkblank[1]{\ifthenelse{\DTLiseq{#1}{}}{***}{#1}}
\makeatother

\begin{document}
	\linespread{2.5}
	\DTLloaddb{addresses}{db/questions.csv} % use your actual address database here
	\DTLforeach*{addresses}{%
		\VarOOA  =VarOOA,
		\VarOOB  =VarOOB,
		\VarOOC  =VarOOC,
		\VarOOD  =VarOOD,
		\VarOOE  =VarOOE,
		\VarOOF  =VarOOF,
		\VarOOG  =VarOOG,
		\VarOOH  =VarOOH,
		\VarOOI  =VarOOI,
		\VarOOJ  =VarOOJ,
		\VarOOK  =VarOOK,
		\VarOOL  =VarOOL,
		\VarOOM  =VarOOM,
		\VarOON  =VarOON,
		\VarOOO  =VarOOO,
		\VarOOP  =VarOOP,
		\VarOOQ  =VarOOQ,
		\VarOOR  =VarOOR,
		\VarOOS  =VarOOS,
		\VarOOT  =VarOOT,
		\VarOOU  =VarOOU,
		\VarOOV  =VarOOV,
		\VarOOW  =VarOOW,
		\VarOOX  =VarOOX,
		\VarOOY  =VarOOY,
		\VarOOZ  =VarOOZ,
		\VarOOa  =VarOOa,
		\VarOOb  =VarOOb,
		\VarOOc  =VarOOc,
		\VarOOd  =VarOOd,
		\VarOOe  =VarOOe,
		\VarOOf  =VarOOf,
		\VarOOg  =VarOOg,
		\VarOOh  =VarOOh,
		\VarOOi  =VarOOi,
		\VarOOj  =VarOOj,
		\VarOOk  =VarOOk,
		\VarOOl  =VarOOl,
		\VarOOm  =VarOOm,
		\VarOOn  =VarOOn,
		\VarOOo  =VarOOo,
		\VarOOp  =VarOOp,
		\VarOOq  =VarOOq,
		\VarOOr  =VarOOr,
		\VarOOs  =VarOOs,
		\VarOOt  =VarOOt,
		\VarOOu  =VarOOu,
		\VarOOv  =VarOOv,
		\VarOOw  =VarOOw,
		\VarOOx  =VarOOx,
		\VarOOy  =VarOOy,
		\VarOOz  =VarOOz,
		\SName =Name,
		\No =No}
	{%
		\setcounter{page}{1}
		%\thispagestyle{empty}
		\begin{center}
			\begin{minipage}{.15\textwidth}
				\begin{flushleft}
					\includegraphics[width=\textwidth]{images/uadec-original}
				\end{flushleft}
			\end{minipage}
			\begin{minipage}{.84\textwidth}
				\begin{flushright}
					{\Huge \textbf{Universidad Autónoma de Coahuila}}\\[2mm]
					{\huge Facultad de Ingeniería Mecánica y Eléctrica}\\[2mm]
					{\LARGE Unidad Torreón}
				\end{flushright}
			\end{minipage}
		\end{center}
		%\\[0.5cm]
		\begin{center}
			\setlength\doublerulesep{2pt}\doublerulesepcolor{LightCyan}
			\begin{tabularx}{\textwidth}{ ||>{\columncolor{Gray}}l|X||>{\columncolor{Gray}}l|r|| }
				\hhline{|t==:t:==t|}
				Subject      		& \Subject  		& Group         	& \Group   					\\ \hhline{|:==::==:|}
				Degree         		& \Carrera  		& Date      		& \Date     				\\ \hhline{|:==::==:|}
				Exam / Homework		& \ExamType    		& Registration \#	& \textbf{\textit{\No}}       				\\ \hhline{|:==::==:|}
				Professor's name	& \PName			& Marks Obtained	& \underline{\hspace{1cm}} $\Big /10$				\\ \hhline{|:==:b:==:|}
				Student's name		& \multicolumn{3}{X||}{\textbf{\textit{\MakeUppercase{\SName}}}}	\\ \hhline{|b====b|}
			\end{tabularx}
		\end{center}
		\section{Important dates}
		\begin{tabularx}{\textwidth}{|c|l|l|l|X|}
			\hline
			\textbf{Activity \#}	&	\textbf{From}	&	\textbf{To}	&	\textbf{Responsible}	&	\textbf{Activity}	\\ \hline
			1	&	27-Mar	&	28-Mar	&	Editor / professor	&	Group formation and assignment of the chapters.	\\ \hline
			2	&	28-Mar	&	07-Apr	&	Authors / student	&	Writing the chapter contents.	\\ \hline
			3	&	07-Apr	&	28-Apr	&	Reviewers / student	&	Review the chapter contents.	\\ \hline
			4	&	28-Apr	&	05-May	&	Authors / student	&	Author should make the corrections given by the reviewers.	\\ \hline
			5	&	05-May	&	08-May	&	Editor / professor	&	Review the book and give corrections to the authors. 	\\ \hline
			6	&	08-May	&	11-May	&	Authors / student	&	Author should make the corrections suggested by editor.	\\ \hline
			7	&	11-May	&	15-May	&	Editor / professor	&	Final version is prepared and published on the website. 	\\ \hline
		\end{tabularx}
		\section{Chapters}
		\begin{enumerate}
			\item Semiconductor materials
			\item Diodes
			\item Bipolar Junction Transistor
			\item Field Effect Transistor
			\item Experiment design to characterize semiconductor components
		\end{enumerate}
		\section{Assignation}
			You are assigned to the Group {\VarOOA}, which is author to the Chapter {\VarOOA} and reviewer to the chapter {\VarOOB}.
		\section{Instructions}
		\begin{enumerate}
			\item It is duty of the group members to divide work among themselves for writing and reviewing.
			\item You should write in Microsoft Word. (A free copy of Microsoft office is available to the students of UAdeC.)
			\item Text of each chapter should at-least contain 15 pages. ( including the figures and tables, the chapter may be longer than 15 pages)
			\item The page should use a margin of 1 inch on all the sides.
			\item Use the font "Times New Roman", size 11 for normal text. The font size for the first heading should be 28, second heading should be 18.
			\item \textbf{Late submission} will not be tolerated, it will lead to \textbf{zero marks}.
			\item Authors and reviewers should exchange documents only through emails with a copy to the professor "gadis-at-uadec-dot-edu-dot-mx".
			\item Work would be seen as a group task, any mistake in a group will effect all he members of the group.
		\end{enumerate}
		\section{Rating rubric}
		\begin{enumerate}
			\item This work contain a sum of 100 marks.
			\item 25 marks for writing.
			\item 25 marks for reviewing.
			\item 25 marks for the presentation (i.e. following all the instructions of fonts, writing caption to the figures and tables etc.).
			\item 25 marks for the corrections.
			\item Marks will be deducted for the following.
			\begin{enumerate}
				\item Any kind of \textbf{plagiarism} in the text or images will be punished. Reviewers will be awarded \textbf{Zero} marks and the authors will be awarded 50\% of marks.
				\item Sending document after the deadline will result in zero marks.
				\item Each incorrect statement, explanation, text or image will reduce 5 marks to the reviewer and 5 marks to the author.
				\item Each spelling mistake will reduce 10 marks to reviewer and 10 marks to the author.
			\end{enumerate}
		\end{enumerate}
				
		
		\clearpage
	}
\end{document}
